\section{Coverage and source audit}\label{audit}
\subsection{Scope and provenance}
The six weekly PDFs contain 54 pages. The main requested textbook boundary is section 12.2 (printed pp.~702--711; PDF pp.~722--731) and section 12.3 (printed pp.~711--725; PDF pp.~731--745) of the 8th-edition text. Textbook content after the 12.4 heading on printed p.~725 is excluded. The supplementary preview is drawn from weekly lecture files only. The user supplied 13 suggested exercises from each section; all 26 appear in \hyperref[textbook]{textbook practice}. Graphs in \texttt{graphs/01--10} are original labelled reconstructions generated by \texttt{graphs/generate\_graphs.py}; \texttt{graphs/source-*} are cropped source diagrams from the textbook. The diagrams serve the questions shown beside them. All mathematical content has been checked independently against page images and text extraction.

\subsection{Page-to-topic checklist}
\small
\begin{longtable}{p{.40\linewidth}p{.50\linewidth}}
\toprule Source and page(s) & Covered at\\\midrule\endhead
\texttt{Week2.pdf} p.~1 & \hyperref[surface]{plane surface}, \hyperref[cross]{cross-sections}, weekly page 1\\
\texttt{Week2.pdf} p.~2 & \hyperref[surface]{surface test}, \hyperref[cylinder]{cylinders}, weekly page 2\\
\texttt{Week2.pdf} p.~3 & \hyperref[cross]{six trace family}, weekly page 3\\
\texttt{Week2.pdf} pp.~4--5 & \hyperref[contour]{definitions}, \hyperref[radial]{bowl/cone levels}, weekly pages 4--5\\
\texttt{Week2.pdf} p.~6 & \hyperref[plane]{parallel plane contours}, weekly page 6\\
\texttt{Week2.pdf} pp.~7--8 & \hyperref[saddle]{complete table and saddle contours}, weekly pages 7--8\\
\texttt{Week2.pdf} pp.~9--10 & \hyperref[tent]{finite tent roof}, weekly pages 9--10\\
\texttt{sep+14.pdf} pp.~1--4 & \hyperref[graph]{graph test}, \hyperref[transform]{transformation}, \hyperref[cross]{saddle slices}, weekly source section\\
\texttt{sep+16.pdf} pp.~1--4 & \hyperref[cylinder]{free-variable cylinder}, \hyperref[radial]{downward bowl contours}, weekly source section\\
\texttt{sep+18.pdf} pp.~1--4 & \hyperref[table]{full numeric interpolation table}, \hyperref[saddle]{saddle comparison}, weekly source section\\
\texttt{MAT235H-5201-Sept-16.pdf} pp.~1--6 & \hyperref[supplement]{12.1 recap} and weekly review questions; p.~1 administrative information only\\
Same PDF pp.~7--12 & Complete wind-chill table, coordinate questions, heater cross-sections, plane in \hyperref[weekly]{weekly questions}; \hyperref[cross]{method}\\
Same PDF pp.~13--19 & \hyperref[cylinder]{cylinders}, \hyperref[cross]{traces and upper-hemisphere slice}, weekly discussion\\
Same PDF pp.~20--23 & \hyperref[radial]{contours}, \hyperref[plane]{linear levels}, \hyperref[saddle]{table and saddle}, weekly questions\\
Same PDF pp.~24--27 & \hyperref[preview]{12.4 preview} and \hyperref[weekly]{airline/plane/table/surface questions}\\
\texttt{MAT235H-5201-Sept-17.pdf} pp.~1--3 & \hyperref[preview]{plane/cylinder/sphere preview}, weekly questions\\
Same PDF pp.~4--5 & \hyperref[preview]{12.5 three-input levels}, weekly question and self-test\\
Textbook 12.2 pp.~702--707 & \hyperref[graph]{graphs}, \hyperref[transform]{shifts and bell}, \hyperref[cross]{saddle and plane traces}, \hyperref[cylinder]{cylinder example}\\
Textbook 12.2 pp.~707--711 & All user-listed odd exercises in \hyperref[textbook]{practice}; images where questions require them\\
Textbook 12.3 pp.~711--718 & \hyperref[estimate]{corn map/estimation}, \hyperref[radial]{cone/bowl}, \hyperref[saddle]{saddle table}, \hyperref[supplement]{production}\\
Textbook 12.3 pp.~718--725 & All user-listed odd exercises in \hyperref[textbook]{practice}; images where questions require them\\
\bottomrule\end{longtable}\normalsize

\subsection{Problem-type and exercise checklist}
Every \emph{Problem type} section in this file has a source-labelled complete question, worked solution, separately labelled answer, relevant diagram/table if needed, and a common-mistake note. The problem types are \hyperref[surface]{surface description}, sphere/hemisphere test, \hyperref[transform]{transformation}, parabolic traces, nonlinear trace families, \hyperref[cylinder]{cylinder identification}, cross-section vs contour, \hyperref[radial]{algebraic radial contours}, \hyperref[plane]{linear contours}, \hyperref[estimate]{map estimation}, \hyperref[table]{table-derived contours}, \hyperref[saddle]{saddle table and map}, and \hyperref[tent]{finite roof/distance}. Suggested textbook items are \textbf{12.2:} 1, 3, 5, 7, 9, 11, 13, 15, 17, 21, 27, 31, 41; \textbf{12.3:} 1, 3, 5, 7, 9, 11, 13, 15, 17, 19, 25, 27, 47. Each numbered practice item is checked for all subparts, required source visual, concise work and final answer; it refers back to the corresponding knowledge section.

\subsection{Corrections, ambiguity, and limits}
\begin{itemize}[leftmargin=*]
\item The handwritten \texttt{sep+14.pdf} p.~4 names a $y=a$ slice of a paraboloid a ``paraboloid.'' Its slice is a 2D \emph{parabola} in the plane $y=a$; the full graph is a paraboloid. The fixed-plane condition is retained explicitly here.
\item In the September 16 lecture, p.~17, separate bounds $-1\leq x,y\leq1$ do not define the full domain of the upper unit hemisphere. The exact domain is $x^2+y^2\leq1$.
\item The \texttt{sep+18.pdf} table is consistent with $f=12-2x-y$, but finite samples justify an interpolated contour, not a proven global formula. Answers here say ``approximately'' when interpolation is used.
\item The September 17 first page shows two plane equations without its original verbal prompt. Their graph and intercepts are given as an explicit interpretation, with missing wording flagged rather than invented.
\item No material handwritten formula or diagram remained unreadable after high-resolution inspection. A few crossed-out scribbles carry no mathematical content. Corn-map and other source-figure readings are estimates unless an exact labelled line passes through the point.
\end{itemize}
\subsection{Graphic and build checks}
The generated files \texttt{01-surfaces.png} through \texttt{10-self-test.png} display the core and preview objects, with axis labels and relevant heights. Most 3D plots show finite windows of infinite surfaces, not bounded domains. The tent roof is genuinely finite. The original textbook crops used by practice exercises remain in \texttt{graphs/}. This source is intended to compile with \texttt{pdflatex} from the \texttt{W2} directory; re-run the build after editing figures or text.
